\section{Industry and company analysis}

\begin{qfmethod}[Industry and competitive analysis: defining the industry, classification systems, size, growth, profitability, market share and concentration, the five forces and PESTLE, and the industry questions in a valuation]\label{meth:industry-analysis}
\qffield{Statement} \emph{Process.} Porter's typical steps \citep[p.~15]{five-competitive-forces-2008}:
\begin{enumerate}
\item define the relevant industry by product scope and geographic scope;
\item identify buyers, suppliers, competitors, substitutes and potential entrants, segmenting them where useful;
\item assess the drivers of each force;
\item determine overall structure and test it against actual long-run profitability;
\item analyse likely changes in each force;
\item identify what competitors, entrants or the company could change.
\end{enumerate}
Defining the industry too broadly hides differences that matter to competition. Defining it too narrowly misses linkages across related products and regions. If two products share buyers, suppliers and barriers, they belong to one industry (p.~14).

\emph{Classification.} Analysts group companies to compare them.
\begin{itemize}
\item \emph{Statistical} approaches cluster companies whose returns are correlated. They are backward-looking and can produce unintuitive groups \citep[p.~13]{gics-methodology-2023}.
\item \emph{Production-oriented} systems group by what companies make.
\item \emph{Market-oriented} systems group by what customers buy.
\end{itemize}
GICS is demand-oriented, with 11 sectors, 25 industry groups, 74 industries and 163 sub-industries (pp.~5--6, 13). A company is classified by the sub-industry producing more than 60 percent of its revenue, with earnings and market perception as secondary tests (pp.~15--16). A power-tool maker falls in Household Appliances (25201040), which includes manufacturers of power and hand tools, in Consumer Discretionary (p.~29). A maker of industrial machine tools falls in Industrial Machinery (20106020, p.~24). Commercial classifications differ, and peers for valuation may need to cross them (\hyperref[meth:peer-group-selection]{peer groups}).

\emph{Size, growth, profitability and share.} Industry size is the sum of participants' sales in the defined market, and growth is its compound rate. Profitability is the distribution of margins and returns on capital across participants. A company's market share is its sales over industry sales. Concentration is measured by concentration ratios or the Herfindahl--Hirschman index, the sum of squared shares in percent. US agencies treat markets above 1,800 as highly concentrated and changes above 100 as significant; a merger of firms with shares $a$ and $b$ raises the index by $2ab$ \citep[pp.~5--6 and nn.~13--14]{merger-guidelines-2023}.

\emph{Five forces.} Industry profitability in the medium and long run depends on five forces \citep[pp.~2--3, 10]{five-competitive-forces-2008}:
\begin{itemize}
\item the threat of entry, which caps prices or forces investment (p.~3);
\item the power of suppliers, who capture value where cost increases cannot be passed on (p.~6);
\item the power of buyers, who force prices down or service up (p.~7);
\item the threat of substitutes (p.~8);
\item rivalry among existing competitors (p.~9).
\end{itemize}
Growth is not the same as attractiveness, and government is not a sixth force but acts through the five (p.~10).

\emph{PESTLE} scans the external environment for political, economic, social, technological, legal and environmental factors. The CIPD's nine steps run from setting the scope and choosing sources, through gathering and analysing information and ranking factors by importance and risk, to identifying options, sharing the findings and monitoring actions. Its limits: the scan can oversimplify, become an exercise in itself, rest on assumptions, and go out of date quickly \citep[sections on the definition, the nine steps and the limitations]{cipd-pestle-2026}.

\emph{Questions for a valuation} (Interpretation): how big and fast-growing is the market, and for how long; what returns can participants earn, and are they above the cost of capital; which forces will compress them; and how exposed is the industry to the business cycle, technology and regulation. The answers set the long-run growth rate, the margin the company converges to and the length of its competitive advantage period (\hyperref[meth:forecast-horizon-steady-state]{forecast horizon}).

\qffield{Practice} \emph{Original: European power tools}, sales in EUR millions. Harbour has 1,080, Arden 1,500, Brook 900, Corran 700, and many small firms 2,820.
\begin{itemize}
\item \emph{Size and growth.} The industry is 7,000, up from 6,000 four years earlier, 3.93 percent a year.
\item \emph{Shares.} Harbour has 15.43 percent. The HHI of the four named firms is 963, so the market is not concentrated.
\item \emph{A merger.} Harbour combining with Corran would add $2\times15.43\times10.00=309$ points: significant, but the post-merger index is far below 1,800.
\end{itemize}

Five forces:
\begin{itemize}
\item \emph{Buyers} are strong: large retailers and distributors consolidate purchasing.
\item \emph{Suppliers} are moderate: steel, motors and battery cells, with cells concentrated.
\item \emph{Entry} is moderate: Asian manufacturers enter through retailers' own brands.
\item \emph{Substitutes} are weak for trade users and stronger for home users, who can rent.
\item \emph{Rivalry} is high in the home segment, which is price-led, and lower in the trade segment, which is led by battery platforms.
\end{itemize}
Long-run margins should therefore converge to modest levels in home tools and higher in trade tools.

PESTLE flags:
\begin{itemize}
\item tariffs on imported components (political);
\item construction cycles (economic);
\item do-it-yourself trends (social);
\item battery technology (technological);
\item product-safety and battery-collection rules (legal);
\item energy and carbon costs (environmental).
\end{itemize}

\qffield{Conditions} Concentration thresholds are the US agencies' December 2023 guidelines, a framework for merger review, not a measure of industry profitability. The Porter reprint is cited by article page. Industry data in the example are illustrative.

\qffield{Pitfalls} Using a classification code as the industry definition. Treating high growth as high attractiveness. Listing PESTLE factors without ranking them or linking them to cash flows. Computing market share against a market defined to flatter the company.

\qffield{See also} \hyperref[meth:company-position-analysis]{company position}; \hyperref[meth:competition-inflation-forecasting]{competition and inflation in forecasts}; \hyperref[def:business-model-features-types]{business models}.

\end{qfmethod}

\begin{qfmethod}[Company analysis: competitive strategy and position, revenue drivers, profitability and pricing power, and evaluating working capital, investment and capital structure]\label{meth:company-position-analysis}
\qffield{Statement} \emph{Strategy and position.} Strategy builds defences against the competitive forces, or finds a position where they are weakest \citep[p.~12]{five-competitive-forces-2008}. In heavy trucks, where large fleet buyers push prices down, Paccar focused on owner-operators, who have little buying power and are less price-sensitive. It built customised trucks with features they value, and they paid about a 10 percent premium (pp.~12--13). Companies can also exploit industry change and reshape structure (pp.~13--14). Positions are commonly described as cost leadership, differentiation or focus (Interpretation). The test of a position is whether it earns returns above the cost of capital that rivals cannot copy.

\emph{Revenue drivers.} Revenue is decomposed into volume and price, by segment and geography (\hyperref[ex:harbour-pro-forma-model]{pro forma model}). Its drivers are market growth, market share and price. Pricing power is the ability to pass cost increases to customers without losing volume. It shows up in stable gross margins when input costs rise (\hyperref[meth:competition-inflation-forecasting]{pass-through}).

\emph{Profitability.} Analyse margins and the return on invested capital, decomposed into margin and capital turnover, and compare with peers and the cost of capital (\hyperref[meth:return-on-equity-decomposition]{DuPont}). Distinguish cyclical from structural changes.

\emph{Working capital, investment and capital structure.}
\begin{itemize}
\item \emph{Working capital}: the cash cycle against peers and the trend (\hyperref[meth:cash-conversion-cycle]{cash cycle}).
\item \emph{Investment}: capital expenditure relative to depreciation and to sales growth, the return on new capital, and acquisitions (\hyperref[meth:capital-project-evaluation]{projects}).
\item \emph{Capital structure}: leverage, coverage and maturity against cash-flow stability and covenants (\hyperref[meth:capital-structure-benchmarking]{benchmarking}).
\end{itemize}
Each of these feeds a forecast assumption (\hyperref[meth:ratio-based-earnings-forecast]{ratio-based forecasts}).

\qffield{Practice} \emph{Harbour Tools} (\hyperref[ex:harbour-tools-statements]{case}).
\begin{itemize}
\item \emph{Position.} A mid-sized maker with 15 percent of its market, strong in trade tools with a battery platform and weaker in price-led home tools.
\item \emph{Revenue.} 2025 revenue grew 8 percent to 1,080. The forecast splits volume and price at 3 percent each.
\item \emph{Pricing power.} In the case's cost-shock analysis, a low pass-through compresses gross margin. Strong retail buyers limit pricing power in the home segment.
\item \emph{Profitability.} Operating margin fell from 11.0 to 8.6 percent, 10.3 percent before the restructuring charge. Return on net operating assets fell from 8.53 to 6.85 percent, below the 8 percent WACC.
\item \emph{Working capital.} A 75-day cash cycle against a peer's 55 ties up capital, mostly in inventory.
\item \emph{Investment.} 2025 capital spending of 127 was well above depreciation; the bolt-on acquisition adds goodwill and integration risk.
\item \emph{Capital structure.} Net debt to EBITDA of 1.63, rising to 1.85 with the acquisition, within a 2.50 covenant (\hyperref[ex:integrated-credit-decision]{credit decision}).
\end{itemize}
Conclusion: Harbour's returns are below its cost of capital, and value depends on margin recovery and on investment earning more than in 2025. The forecast should assume convergence towards peer returns, not beyond them, unless the trade segment's position strengthens.

\qffield{Conditions} The figures are from the hypothetical case. A position described in words should be checked against numbers: returns, margins through a cycle and market-share trends.

\qffield{Pitfalls} Describing strategy without evidence in returns. Forecasting margin expansion without a source of pricing power or cost advantage. Ignoring working capital and investment when growth is forecast.

\qffield{See also} \hyperref[meth:industry-analysis]{industry analysis}; \hyperref[def:business-model-features-types]{business models}; \hyperref[meth:forecast-horizon-steady-state]{forecast horizon}.

\end{qfmethod}
